\documentclass[conference]{IEEEtran}
\IEEEoverridecommandlockouts

\usepackage{amsmath,amssymb}
\usepackage{booktabs}
\usepackage{cite}
\usepackage{graphicx}
\usepackage{url}
\usepackage{xcolor}

\newcommand{\todo}[1]{\textcolor{red}{[TBD]}}
\newcommand{\thz}{THZ1}
\newcommand{\thp}{THP1}

\title{Event-Driven Reinforcement Learning for Shared Demand-Budget Charging of Multi-Line Catenary-Free Tramways}

\author{
\IEEEauthorblockN{Jun-Jie Zhao}
\IEEEauthorblockA{School of Electrical and Computer Engineering\\
The University of Sydney\\
Sydney, Australia\\
Email: \todo{author email}}
}

\begin{document}
\maketitle

\begin{abstract}
Catenary-free tramways shift traction-energy management from wayside supply to repeated station charging, making dwell-time and charging-power decisions operationally coupled with passenger service and peak-demand exposure. This paper studies a multi-line charging-control problem motivated by two Guangzhou light-rail tramway systems: Haizhu Tram Line 1, a supercapacitor-only line, and Huangpu Tram Line 1, a hybrid supercapacitor--lithium-titanate-battery line. We build an event-driven discrete-event simulation environment in which each tram arrival triggers a continuous control action over dwell time and total station-charging power. The two lines are mixed on a common global event timeline and coupled only through a shared aggregate demand-budget penalty, without assuming common electrical infrastructure or performing power-flow simulation. The proposed formulation extends a single-line event-driven Markov decision process to a 13-dimensional multi-line state that exposes cross-line and own-line charging context while preserving a single-vehicle action space. The evaluation protocol covers a proximal policy optimization controller, fixed-operation control, genetic algorithm, particle swarm optimization, differential evolution, and model predictive control baselines under four running-time disturbance settings and shared-versus-split budget accounting. Experiments are currently under execution; this manuscript reports the formulation, data construction protocol, evaluation design, and result placeholders for an AUPEC conference submission.
\end{abstract}

\begin{IEEEkeywords}
Catenary-free tramway, electric public transport, demand management, reinforcement learning, discrete-event simulation, proximal policy optimization.
\end{IEEEkeywords}

\section{Introduction}
Catenary-free light rail and tramway systems reduce visual intrusion and fixed electrification assets along urban corridors, but they make vehicle energy management more operationally active. Vehicles must recover sufficient energy during station stops, while the same stops also serve passenger boarding and alighting. A controller that increases charging power or dwell time may improve energy feasibility, but it can increase peak aggregate demand, disrupt headways, and affect passenger waiting. These trade-offs are especially relevant for operators that manage multiple electrified tram corridors under a common demand-budget or electricity-procurement policy.

Most charging-control studies for railway or transit electrification treat a single line, a fixed timetable, or a predetermined electrical network model. This paper instead targets a newly deployed engineering scenario: two catenary-free tramway lines with different onboard storage flexibility are operated on the same simulated morning peak and evaluated under a shared aggregate charging-power budget. The setting is motivated by Guangzhou Haizhu Tram Line 1 (\thz), which uses station charging with onboard supercapacitors, and Guangzhou Huangpu Tram Line 1 (\thp), which uses hybrid supercapacitor and lithium-titanate battery storage \cite{haizhu_wiki,haizhu_baike,huangpu_wiki,huangpu_sasac}. The shared budget is a research setting for unified demand management; it is not a claim that the two geographically separated lines share substations, feeders, or physical electrical infrastructure.

The central technical challenge is that charging decisions are event-driven and asynchronous. At each arrival event, the controller sees the currently active charging intervals and the current vehicle state, then chooses only the dwell time and total charging power for the arriving vehicle. The resulting aggregate power depends on overlapping charging intervals from both lines. Enforcing a hard aggregate cap would require an additional multi-vehicle reallocation rule that is not part of the event action. We therefore model the demand budget as a soft penalty on post-action aggregate power, creating a peak-avoidance training signal while preserving a single-event decision structure.

This paper makes three contributions. First, it formulates a two-line event-driven Markov decision process (MDP) for catenary-free tramway charging, extending the single-line state to include total aggregate power, own-line aggregate power, active charging count, remaining charging demand, and line identity. Second, it develops a reproducible discrete-event simulation environment using public line information, explicit engineering assumptions, passenger-demand profiles, traction-energy dynamics, and hybrid-storage energy-management rules. Third, it defines an AUPEC-style evaluation protocol comparing PPO with fixed-operation, GA, PSO, DE, and MPC baselines across disturbance environments and shared-versus-split demand-budget accounting. The numerical results are intentionally left as TBD until the current experiments complete.

\section{Related Work}
Charging and energy-management studies for electric rail and transit systems commonly optimize timetable adherence, energy consumption, regenerative-energy use, or charging schedules under infrastructure constraints. In catenary-free tramways, the key operational distinction is that station dwell time and charging power are coupled: each charging decision also affects passenger service and subsequent arrival times. This motivates an event-driven formulation rather than a purely time-indexed charging schedule.

Metaheuristic optimizers such as GA, PSO, and DE provide natural baselines for continuous dwell-time and charging-power decisions because they can optimize a full-episode action vector without differentiating through the simulator \cite{holland1975adaptation,kennedy1995pso,storn1997de}. Their limitation in the present setting is that a complete candidate vector is tied to a reference episode structure and must be re-optimized when the accounting rule or scenario changes. MPC provides a stronger online optimization baseline by repeatedly solving a finite-horizon problem, but it uses explicit future-event information. PPO is used as the learning baseline because it supports continuous stochastic policies and can condition decisions on the current event state \cite{schulman2017ppo}.

The present work differs from these generic control choices by focusing on cross-line demand-budget coupling without extending the action to a centralized multi-vehicle dispatch. The novelty is therefore in the event-driven two-line formulation and evaluation setting, not in proposing a new RL algorithm.

\section{System Overview}
\subsection{Two-Line Catenary-Free Tramway Setting}
The simulated system contains two line containers. The \thz{} case represents Haizhu Tram Line 1 with an 11-station, 7.7 km opening-period alignment, seven vehicles, and supercapacitor-only onboard storage. The \thp{} case represents Huangpu Tram Line 1 with 19 stations, approximately 14.3 km route length, 11 peak-period vehicles, and hybrid supercapacitor--lithium-titanate battery storage. The line properties are derived from public sources where available and otherwise specified as explicit assumptions in the input data files. All station-level passenger rates used by the simulator are assumed morning-peak curves, not measured passenger counts.

Both lines are aligned to a common morning-peak simulation window from 07:30 to 09:30. The simulator initializes a fixed fleet cycle for each line rather than generating vehicles from a fixed headway process. Vehicles are distributed around their line cycle at the start of the window; headways and event overlaps then emerge from fleet size, section running times, dwell choices, turnarounds, and disturbances.

\subsection{Event-Driven Decision Process}
The simulation is a discrete-event system. Each event $k$ corresponds to tram $i$ arriving at station platform $n$ at time $t_k$. The controller outputs
\begin{equation}
    a_k = (d_k, P_k),
\end{equation}
where $d_k$ is dwell time and $P_k$ is total grid-side station-charging power for the arriving vehicle. The method-layer action bounds are
\begin{equation}
    10 \le d_k \le 50 \quad \mathrm{s}, \qquad 0 \le P_k \le 300 \quad \mathrm{kW}.
\end{equation}
The station physical availability and storage feasibility can only reduce the selected power. Stops without charging infrastructure would force $P_k=0$; in the current two-line case all stations are assumed charging-capable. A small dead zone maps powers below 5 kW to exactly zero in the two-line environment, allowing explicit skip-charging actions.

Events from both lines are inserted into one global priority queue. Events are processed by increasing arrival time, with deterministic tie-breaking by line code, station identifier, and tram identifier. Effective charging intervals are half-open intervals $[t_k,t_k+d_k)$, so a charging interval created by an earlier same-time event is immediately visible to later same-time events.

\begin{figure}[t]
\centering
\fbox{\parbox{0.92\columnwidth}{\centering\vspace{1.5em}TBD: event-driven two-line simulation pipeline, showing arrival event, state construction, policy action, EMS projection, charging-interval registration, and shared/split overload accounting.\vspace{1.5em}}}
\caption{Planned system diagram for the two-line event-driven charging-control environment. The final version will replace this placeholder with the simulator pipeline figure.}
\label{fig:pipeline}
\end{figure}

\section{Model Formulation}
\subsection{State Representation}
The two-line MDP uses a 13-dimensional pre-action state:
\begin{align}
 s_k = (&t_k, i, n, h_k, E^{\mathrm{sc}}_{i,k}, E^{\mathrm{bat}}_{i,k}, q_{i,k}, w_{n,k}, 
 P^{\mathrm{tot}}_{k}, P^{\mathrm{own}}_{k}, 
 N^{\mathrm{ch}}_{k}, E^{\mathrm{rem}}_{k}, \ell_k),
\end{align}
where $h_k$ is the headway since the previous same-line, same-station, same-direction departure; $E^{\mathrm{sc}}$ and $E^{\mathrm{bat}}$ are onboard supercapacitor and battery energies; $q$ is onboard passengers; $w$ is waiting passengers; $P^{\mathrm{tot}}$ is total active charging power across both lines; $P^{\mathrm{own}}$ is active charging power on the current line; $N^{\mathrm{ch}}$ is the cross-line active charging count; $E^{\mathrm{rem}}$ is the remaining grid-side planned energy in active charging intervals; and $\ell_k$ is a line code. Vehicle and station identifiers remain local to each line. Direction is intentionally excluded from the state and enters the dynamics through same-platform headway bookkeeping.

For an active interval $r$ with projected power $P_r$, start time $\tau_r$, and dwell $d_r$, the remaining grid-side energy at event time $t$ is
\begin{equation}
 E^{\mathrm{rem}}(t) = \sum_{r \in \mathcal{A}(t)} \max\left(0, \frac{P_r d_r}{3600} - \frac{P_r \min(\max(t-\tau_r,0),d_r)}{3600}\right),
\end{equation}
where $\mathcal{A}(t)$ is the set of active charging intervals. This quantity describes remaining scheduled grid-side demand, not storage headroom.

\subsection{Passenger and Time Dynamics}
For each platform, time-varying boarding and alighting demand rates are integrated between the previous same-platform departure and the current arrival. With service time $\delta$ per passenger, the event service budget is $d_k/\delta$. Alighting is limited by onboard passengers and service budget; boarding is limited by waiting demand, residual vehicle capacity, and remaining service budget. The post-event onboard and waiting counts are then used for passenger-cost and subsequent energy calculations.

The departure and next-arrival times follow
\begin{align}
    t^{\mathrm{dep}}_k &= t_k + d_k, \\
    t^{\mathrm{arr}}_{k+1} &= t^{\mathrm{dep}}_k + T^{\mathrm{run}}_{k},
\end{align}
where the realized running time is either the planned section time or a delayed value under the disturbance model. The four disturbance environments are $E1=(0.2,1.10)$, $E2=(0.5,1.10)$, $E3=(0.2,1.20)$, and $E4=(0.5,1.20)$, where each pair denotes Bernoulli delay probability and multiplicative delay factor. Disturbance random streams are derived independently by line from the global seed.

\subsection{Energy and Storage Dynamics}
The traction-energy model solves a triangular or trapezoidal speed profile for each section and integrates acceleration, cruise, and braking phases in SI units. Acceleration and cruise energy are divided by traction efficiency; regenerative braking is subtracted using a regeneration efficiency. Auxiliary energy is proportional to running time and onboard passenger count:
\begin{equation}
    W^{\mathrm{aux}}_k = \left(P_{\mathrm{ic}} + P_{\mathrm{oc}} q_{i,k}\right)\frac{T^{\mathrm{run}}_k}{3600}.
\end{equation}
The total section demand is $W_k = W^{\mathrm{trac}}_k + W^{\mathrm{aux}}_k$.

For station charging, grid-side energy $P_k d_k/3600$ is converted to stored energy through component-specific charging efficiencies. In the supercapacitor-only case, charging is projected to available supercapacitor headroom. In the hybrid case, a deterministic onboard energy-management system first fills the supercapacitor within its state-of-charge limits and then sends remaining feasible charging energy to the battery subject to battery state-of-charge and C-rate limits. During traction discharge, the supercapacitor is used first and the battery supplies residual demand when necessary; any remaining unmet energy is recorded as emergency deficit. This priority rule is deterministic and introduces no additional policy action dimension.

\subsection{Cost and Reward}
Each event cost contains charging, overcharge, overload, emergency, and passenger-waiting terms:
\begin{equation}
 C_k = C^{\mathrm{ch}}_k + C^{\mathrm{oc}}_k + C^{\mathrm{ol}}_k + C^{\mathrm{em}}_k + C^{\mathrm{wt}}_k,
 \qquad R_k=-C_k.
\end{equation}
The regular charging cost is
\begin{equation}
 C^{\mathrm{ch}}_k = c_e E^{\mathrm{ch}}_k,
\end{equation}
where $E^{\mathrm{ch}}_k$ is stored-side charged energy. Overcharge cost is computed from positive increments of per-vehicle, per-storage-component cumulative excess charging beyond $1.25$ times cumulative consumption. Emergency cost uses only positive post-discharge deficit,
\begin{equation}
 C^{\mathrm{em}}_k = c_m \max(0, W_k - E^{\mathrm{available}}_k - E^{\mathrm{ch,available}}_k),
\end{equation}
so energy-sufficient events do not receive negative emergency cost. Passenger waiting cost is proportional to post-service waiting passengers.

The shared demand-budget cost is evaluated after registering the current event's charging interval:
\begin{equation}
 C^{\mathrm{ol}}_{\mathrm{shared}}(t_k;B) = c_o \max(0, P^{\mathrm{tot}}(t_k)-B)^2.
\end{equation}
The default shared budget is $B=1600$ kW, interpreted as an assumed operator-level demand budget. For split accounting, the same total budget is divided evenly across the two lines:
\begin{equation}
 C^{\mathrm{ol}}_{\mathrm{split}}(t_k;B) = \sum_{m \in \{\thz,\thp\}} c_o \max(0, P_m(t_k)-B/2)^2.
\end{equation}
Both equations are soft penalties, not hard power caps. The default cost weights are $c_e=5$, $c_o=10$, $c_m=1000$, $c_p=750$, and $c_w=20$. PPO uses event-index discount $\gamma=0.993$; search baselines use undiscounted total reward as their optimization objective.

\section{Control Methods and Baselines}
\subsection{PPO Controller}
The learning controller is proximal policy optimization (PPO) with a continuous two-dimensional action. The actor produces an unconstrained latent action, applies a hyperbolic tangent transform, and maps it affinely to the dwell-time and power bounds. Deterministic evaluation uses the actor mean. The actor and critic use three hidden layers of width 512 with Tanh activations, Adam optimization, clipping parameter $\epsilon=0.2$, generalized advantage estimation with $\lambda=0.98$, and initial entropy coefficient $10^{-2}$ \cite{schulman2017ppo}. The learning rate and rollout settings are recorded in run metadata for each experiment.

\subsection{Baselines}
The fixed-operation baseline, denoted w/o-Opt., applies $d_k=30$ s and $P_k=300$ kW at every event before environment-side projection. Three population-based baselines optimize a full-episode candidate vector $x=[d_1,\ldots,d_K,P_1,\ldots,P_K]$ under the same action bounds. Candidate entries are bound to structural keys $(\ell,i,v)$ consisting of line code, tram identifier, and visit index, rather than dynamic arrival order, so the same gene refers to the same physical decision across disturbance realizations. GA uses crossover rate 0.9 and mutation rate 0.1. PSO uses cognitive coefficient 0.5, social coefficient 0.5, and inertia 0.8. DE uses scaling factor 0.5 and recombination rate 0.9 \cite{holland1975adaptation,kennedy1995pso,storn1997de}.

The MPC baseline solves a finite-horizon event optimization repeatedly and applies only the first action. It uses the same simulator, action bounds, cost components, and disturbance profile as the other controllers. Because MPC uses an explicit future-event table, its information advantage is reported separately from learned policy evaluation.

\section{Case Study and Evaluation Design}
\subsection{Input Data and Source Discipline}
The case-study data are divided into public facts, derived quantities, and explicit assumptions. Public facts include line identity, route length, station count, fleet size, vehicle capacity, and storage technology where available. Derived quantities include evenly distributed station distances when exact station-chainage data are unavailable, section running times inferred from route length and published average speed, and supercapacitor energy from capacitance and voltage assumptions. Assumed quantities include station-level passenger-rate profiles, some storage limits, auxiliary power, and the shared demand-budget values. No numerical result from any prior paper is used as a calibration target or validation threshold.

The \thz{} route uses 11 stations and 7.7 km, with seven vehicles and a derived usable supercapacitor capacity of 4.734 kWh. The \thp{} route uses 19 stations and 14.3 km, with 11 peak-period vehicles, supercapacitor capacity of approximately 4.66 kWh, and an assumed 8.0 kWh lithium-titanate battery usable model. Both lines use all-stop charging availability with a method-layer maximum power of 300 kW.

\subsection{Experimental Factors}
The planned experiments vary four dimensions:
\begin{itemize}
    \item controller: w/o-Opt., GA, PSO, DE, MPC, and PPO;
    \item disturbance environment: none, E1, E2, E3, and E4, with E1--E4 used for robustness reporting;
    \item budget value: $B \in \{1600,1400,1200\}$ kW;
    \item accounting rule: shared aggregate budget versus split per-line budget.
\end{itemize}
Each $B$--accounting combination is trained or optimized independently. A policy trained under shared accounting is not reused as a split-accounting result.

\subsection{Metrics and Figures}
For each trained or optimized configuration, evaluation uses seeds $\{0,1,2,3,4\}$. The main event-level metrics are total cost, average event cost, average reward, and the five cost components. Time-series metrics are computed on a 1 s grid over the power-reporting window and include maximum total aggregate power, mean aggregate power, overload duration, line-level power peaks, and line-level overload duration. Event-cost metrics and time-sampled power metrics are reported separately.

The conference paper will contain the following result artifacts once experiments complete:
\begin{itemize}
    \item PPO learning curves under E1--E4 and budget configurations;
    \item aggregate power traces showing $P_{\thz}$, $P_{\thp}$, and $P_{\mathrm{tot}}$ with budget reference lines;
    \item budget-pooling trade-off curves comparing shared and split accounting;
    \item cost-component tables for all controllers;
    \item qualitative flexible-yielding diagnostics, such as \thp{} zero-power decisions during \thz{} charging peaks.
\end{itemize}

\begin{table}[t]
\centering
\caption{Planned Event-Level Cost Summary (Values TBD)}
\label{tab:cost-summary}
\begin{tabular}{llrrrrrr}
\toprule
Method & Env. & $C^{\mathrm{ch}}$ & $C^{\mathrm{oc}}$ & $C^{\mathrm{ol}}$ & $C^{\mathrm{em}}$ & $C^{\mathrm{wt}}$ & Total \\
\midrule
w/o-Opt. & E1--E4 & \todo{} & \todo{} & \todo{} & \todo{} & \todo{} & \todo{} \\
GA & E1--E4 & \todo{} & \todo{} & \todo{} & \todo{} & \todo{} & \todo{} \\
PSO & E1--E4 & \todo{} & \todo{} & \todo{} & \todo{} & \todo{} & \todo{} \\
DE & E1--E4 & \todo{} & \todo{} & \todo{} & \todo{} & \todo{} & \todo{} \\
MPC & E1--E4 & \todo{} & \todo{} & \todo{} & \todo{} & \todo{} & \todo{} \\
PPO & E1--E4 & \todo{} & \todo{} & \todo{} & \todo{} & \todo{} & \todo{} \\
\bottomrule
\end{tabular}
\end{table}

\begin{table}[t]
\centering
\caption{Planned Budget-Pooling Comparison (Values TBD)}
\label{tab:budget-pooling}
\begin{tabular}{llrrrr}
\toprule
$B$ (kW) & Accounting & Peak $P_{\mathrm{tot}}$ & Overload time & Avg. cost & Waiting \\
\midrule
1600 & Shared & \todo{} & \todo{} & \todo{} & \todo{} \\
1600 & Split & \todo{} & \todo{} & \todo{} & \todo{} \\
1400 & Shared & \todo{} & \todo{} & \todo{} & \todo{} \\
1400 & Split & \todo{} & \todo{} & \todo{} & \todo{} \\
1200 & Shared & \todo{} & \todo{} & \todo{} & \todo{} \\
1200 & Split & \todo{} & \todo{} & \todo{} & \todo{} \\
\bottomrule
\end{tabular}
\end{table}

\section{Discussion}
The formulation is deliberately limited to operational demand-budget coordination. It does not model feeder topology, voltage constraints, power-flow equations, depot sharing, cross-line passenger transfer, or vehicle interchange between lines. This scope is appropriate for testing whether a unified operator-level budget can induce cross-line charging coordination, but it cannot support claims about electrical feasibility at substations or distribution feeders. Any future electrical-network conclusion would require a separate network model and data that are not present in the current simulator.

A second limitation is that passenger-rate profiles and several storage parameters are assumed rather than measured. The resulting experiments should therefore be interpreted as simulation studies over a transparent engineering scenario, not as calibrated operational forecasts for Guangzhou. The main validity criterion is internal consistency, reproducibility under fixed seeds, and sensitivity across disturbance and budget settings.

\section{Conclusion}
This paper presents an event-driven simulation and reinforcement-learning formulation for shared demand-budget charging of two catenary-free tramway lines. The method keeps the action space local to a single arriving vehicle while exposing enough cross-line charging context for a shared policy to coordinate aggregate power. The case-study design combines public line information, explicit assumptions, hybrid-storage dynamics, and a soft demand-budget penalty to evaluate PPO and optimization baselines under stochastic running-time disturbances. Experimental results are currently TBD and will be inserted after the planned AUPEC evaluation runs complete.

\bibliographystyle{IEEEtran}
\bibliography{references}

\end{document}
